\documentclass[11pt,a4paper]{article}
\usepackage{amsmath, amssymb, amsfonts}
\usepackage[utf8]{inputenc}
\usepackage[T1]{fontenc}
\usepackage{geometry}
\geometry{margin=2.5cm}
\usepackage{enumitem}
\usepackage{booktabs}
\usepackage{hyperref}
\usepackage{xcolor}

\title{ICNNA Architectural Refactoring\\Phase~2: Architectural Decision --- ICNNA Outcome\\
\large Meeting Minute}
\author{Felipe Orihuela-Espina \and Claude (Computational Research Architect)}
\date{13 April 2026}

\begin{document}
\maketitle

% ================================================================
% Document Log
% ================================================================
\noindent\textbf{Document Log:}
\begin{itemize}[nosep]
    \item 13-Apr-2026: FOE/Claude --- Document created. Records the outcome of Phase~2 as it pertains to ICNNA's future.
\end{itemize}

\vspace{1em}
\hrule
\vspace{1em}

% ================================================================
\section{Session Overview}

\begin{description}[style=nextline]
    \item[Date:] 13 April 2026
    \item[Participants:] Felipe Orihuela-Espina (FOE), Claude Opus 4.6 (Computational Research Architect role)
    \item[Predecessor:] Phase~2a minute (\texttt{claude/Phase2a\_minute.tex}), 12~April~2026.
    \item[Objective:] Record the outcome of Phase~2 as it pertains to ICNNA specifically.
    \item[Outcome:] ICNNA transitions to maintenance mode. Cantor (C++23) becomes the successor project. ICNNA remains live for pressing issues and current users.
    \item[Model used:] Claude Opus 4.6
    \item[Status:] Phase~2 complete for ICNNA. Active development continues in Cantor project.
\end{description}

% ================================================================
\section{Key Decision: Cantor as Successor}

The Phase~2 architectural analysis (Steps~A--C) concluded that:

\begin{enumerate}
    \item \textbf{Option~2 (clean-sheet design)} was selected over evolving ICNNA or a hybrid approach (weighted scores: 204 vs.\ 107.5 vs.\ 149.5).
    \item The implementation language is \textbf{C++23 with wxWidgets}, not MATLAB.
    \item The new system is a \textbf{separate project called Cantor}, not an ICNNA update.
\end{enumerate}

\textbf{Rationale:} ICNNA's value lies in 18~years of domain knowledge and mathematical theory, not in its code or architecture. The mathematical assets (MENA Gen~2/3, pipeline optimization, overprocessing theory) are mostly not in ICNNA's codebase --- they exist in papers and theses. A clean-sheet design in C++23 can be built \emph{from} these mathematical foundations rather than retrofitted to accommodate them.

Full analysis is documented in the Cantor project:
\begin{itemize}
    \item \texttt{Cantor/claude/Phase1\_minute.tex} --- decision framework
    \item \texttt{Cantor/claude/Phase2a\_minute.tex} --- language evaluation
    \item \texttt{Cantor/claude/Phase2b\_minute.tex} --- final decisions and deep analysis
\end{itemize}

% ================================================================
\section{ICNNA's New Role: Maintenance Mode}

ICNNA transitions from active development to \textbf{maintenance mode} with the following scope:

\subsection{In Scope}
\begin{itemize}
    \item \textbf{Bug fixes:} Address bugs reported by current users.
    \item \textbf{Pressing feature requests:} Small, well-scoped features that current users need and that are quick to implement.
    \item \textbf{SNIRF/BIDS compatibility:} Maintain compatibility with current standards so that ICNNA data can be read by Cantor when it is ready.
    \item \textbf{Documentation:} Keep existing documentation current for the features that exist.
    \item \textbf{Version releases:} Continue version numbering (v1.4.x, possibly v1.5.x) for maintenance releases.
\end{itemize}

\subsection{Out of Scope}
\begin{itemize}
    \item Major architectural refactoring.
    \item New analysis methods (these go to Cantor).
    \item GUI overhaul (Cantor will have a new GUI).
    \item Pipeline formalization (FMS, overprocessing --- Cantor scope).
    \item AI integration (Cantor scope).
    \item Performance optimization beyond what is needed for current users.
\end{itemize}

\subsection{Deprecation Timeline}
\begin{itemize}
    \item \textbf{No fixed deprecation date.} ICNNA remains available and functional for as long as users need it.
    \item \textbf{Natural transition:} As Cantor reaches feature parity for basic fNIRS analysis (SNIRF import, preprocessing, connectivity, visualization), users will be guided to transition.
    \item \textbf{ICNNA is not abandoned} --- it is succeeded. All stable releases remain available on GitHub. The repository stays public.
\end{itemize}

% ================================================================
\section{Decisions Recorded}

\begin{table}[h]
\centering
\begin{tabular}{lp{10cm}}
\toprule
\textbf{Item} & \textbf{Decision} \\
\midrule
ICNNA status & Maintenance mode. No major new development. \\
Successor & Cantor (C++23, wxWidgets, MIT license). Separate project. \\
Current users & Supported via ICNNA maintenance releases. \\
SNIRF/BIDS & ICNNA maintains compatibility so data bridges to Cantor. \\
Deprecation & No fixed date. Natural transition as Cantor matures. \\
Repository & Stays public on GitHub. All stable releases preserved. \\
Cantor location & \texttt{C:\textbackslash Users\textbackslash orihuelf-admin\textbackslash OneDrive\textbackslash Git\textbackslash Cantor} \\
\bottomrule
\end{tabular}
\caption{Decisions regarding ICNNA's future.}
\end{table}

% ================================================================
\section{Relationship Between ICNNA and Cantor}

\begin{itemize}
    \item Cantor is \textbf{intellectually descended} from ICNNA but \textbf{architecturally independent}.
    \item No code is shared or ported between the two projects.
    \item Data interoperability is via \textbf{SNIRF/BIDS} standards (file-based exchange).
    \item ICNNA's coding conventions (naming, documentation, logging) inform Cantor's conventions, adapted for C++23.
    \item Domain knowledge from ICNNA (experimental design, fNIRS processing, QC framework) feeds into Cantor's design.
\end{itemize}

% ================================================================
\section{Next Steps for ICNNA}

\begin{enumerate}
    \item Address any pressing issues or bug reports from current users.
    \item Ensure current v1.4.1~beta is stable; consider a v1.4.1 release if warranted.
    \item Maintain SNIRF/BIDS compatibility as standards evolve.
    \item No further architectural planning --- that effort moves to Cantor.
\end{enumerate}

% ================================================================
\section{Documents Referenced}

\begin{enumerate}
    \item Phase~1 minute: \texttt{claude/Phase1\_minute.tex}
    \item Phase~2a minute: \texttt{claude/Phase2a\_minute.tex}
    \item Cantor Phase~2b minute: \texttt{Cantor/claude/Phase2b\_minute.tex}
\end{enumerate}

\end{document}
